\documentclass[10pt,a4paper]{amsart}
\usepackage[margin=19mm]{geometry}
\usepackage{amsmath,amssymb,mathtools}
\usepackage[scr=rsfs,cal=euler]{mathalfa}
\usepackage{xfrac}
\usepackage[unicode,hidelinks,bookmarks=false]{hyperref}
\newcommand{\ie}{i.e.\ }
\newcommand{\cf}{cf.\ }
\newcommand{\resp}{respectively\ }
\newcommand{\Subsection}{\subsection}
\providecommand{\nobreakdash}{\nobreak\hbox{-}}
\setlength{\emergencystretch}{2em}
\multlinegap0pt
\usepackage{bm}
\usepackage{amssymb}
\usepackage{mathtools}
\usepackage[shortlabels]{enumitem}
\newtheorem{theorem}{Theorem}
\newtheorem{lemma}{Lemma}
\newtheorem{proposition}{Proposition}
\newcommand{\GRS}{\operatorname{GRS}}
\newcommand{\N}{\operatorname{N}}
\newcommand{\Tr}{\operatorname{Tr}}
\title{Quantum MDS codes from complements\\of unions of finite-field subsets}
\author{Naihong Hu, Hong Ji}
\date{Source: arXiv:2609.09943v1 (CC BY 4.0)}
\begin{document}
\maketitle

\noindent\emph{Source and licence.} Quantum MDS codes from complements\\of unions of finite-field subsets. Version arXiv:2609.09943v1 (CC BY 4.0), distributed by arXiv under the Creative Commons Attribution 4.0 International licence, http://creativecommons.org/licenses/by/4.0/, as recorded in the arXiv licence metadata for this exact version and shown on the abstract page https://arxiv.org/abs/2609.09943v1. Full licence notice: \texttt{licenses/arxiv-2609.09943-CC-BY-4.0.txt}.\par\smallskip

\noindent\textit{Editorial scope.} Counted unit: the article's theorem thm:fourth-grs (source0.tex lines 638--640). The article's complete original proof of it is delivered with it (source0.tex lines 642--674, one \texttt{proof} environment of 33 lines). No mathematics is written, repaired or re-derived: the statement and the proof are the article's own lines, in the article's own notation, and no retained line is substituted or re-flowed.  Retained uncounted as the source-local context the proof needs: lines 229--239; lines 266--278; lines 373--392; lines 440--451.  Nothing was omitted from any retained span, so the package carries no omission record.  No retained line contains a citation, so the wrapper carries no bibliography and no bibliography entry.  No retained line contains a figure environment, an image-inclusion command or drawing code, so no image is delivered and the package declares no asset dependency.  Editorial formatting only, in addition: an AMS \texttt{amsart} preamble that restates the article's own declarations and macros used by the retained lines, namely the environments \texttt{theorem}, \texttt{lemma}, \texttt{proposition} on separate counters, together with the standard packages those lines need.  The retained environments print in this document as Theorem 1, Lemma 1, Proposition 1 in order of appearance, whereas the article prints the same items with the numbering of its own full document.  The complete native source of the article as distributed in the arXiv e-print for this exact version is bundled unchanged as \texttt{sources/209862/source0.tex}.
\begin{theorem}\label{thm:general-construction}
With the notation above, let $m\ge\varepsilon_X$ be an integer and put $G(x)=x^mP(x)$. Suppose that there exists $\gamma\in U$ such that, for every $\alpha\in A$,
\begin{equation}\label{eq:general-frobenius}
 G(\alpha)\ne0,\qquad G(\alpha)^q=\gamma G(\alpha).
\end{equation}
Then, for every positive integer $k$ satisfying
\begin{equation}\label{eq:general-degree}
 (q+1)k\le q^2+q-1-m-\deg P
\end{equation}
there exists $\boldsymbol v=(v_\alpha)_{\alpha\in A}\in(\mathbb F_{q^2}^{*})^n$ such that $\GRS_k(A,\boldsymbol v)$ is a Hermitian self-orthogonal \mbox{$[n,k,n-k+1]_{q^2}$} MDS code.
\end{theorem}

\begin{lemma}\label{lem:q-power-calculation}
Retain the notation and assumptions of Theorem~\ref{thm:general-construction}, and suppose that $0\notin A$. Let $D\subseteq U$, set $b=|D|\ge1$, and define
\[
 Q(x)=\prod_{\delta\in D}(x^{q-1}-\delta),
 \qquad C_D=(-1)^b\left(\prod_{\delta\in D}\delta\right)^{-1}.
\]
Suppose that $P(x)=H(x)Q(x)$ for some $H\in\mathbb F_{q^2}[x]$ satisfying $H(\alpha)\in\mathbb F_q^{*}$ for every $\alpha\in A$.
Then $C_D\in U$, and for every positive integer $k$ satisfying
\begin{equation}\label{eq:specialized-degree}
 (q+1)k\le q^2+q-1-bq-\deg H
\end{equation}
there exists $\boldsymbol v=(v_\alpha)_{\alpha\in A}\in(\mathbb F_{q^2}^{*})^n$ such that $\GRS_k(A,\boldsymbol v)$ is a Hermitian self-orthogonal \mbox{$[n,k,n-k+1]_{q^2}$} MDS code.
\end{lemma}

\begin{proposition}\label{prop:intersection-formulas}
For all $c,e\in\mathbb F_q$, the sets $L_c$ and $K_e$ each have cardinality $q$ and satisfy
\[
 L_c=\frac c2+\omega\mathbb F_q,\qquad
 K_e=-\frac{e}{2\omega}+\mathbb F_q.
\]
Their intersection consists of the single point
\[
 x_{c,e}=\frac12\left(c-\frac e\omega\right),
 \qquad \N(x_{c,e})=\frac{c^2-e^2/\Delta}{4}.
\]
For every $\delta\in U\setminus\{-1\}$, we have $L_0\cap M_\delta=\varnothing$, and for every $c\in\mathbb F_q^{*}$,
\[
 L_c\cap M_\delta=\left\{\frac{c}{\delta+1}\right\}.
\]
For every $\delta\in U\setminus\{1\}$, we have $K_0\cap M_\delta=\varnothing$, and for every $e\in\mathbb F_q^{*}$,
\[
 K_e\cap M_\delta=\left\{\frac{e}{\omega(\delta-1)}\right\}.
\]
\end{proposition}

\begin{lemma}\label{lem:basic-vanishing-polynomials}
For $c,e\in\mathbb F_q$, $\nu\in\mathbb F_q^{*}$, and $\delta\in U$, the corresponding monic vanishing polynomials are
\begin{equation}\label{eq:basic-vanishing-polynomials}
\begin{aligned}
 F_{L_c}(x)&=x^q+x-c,&
 F_{\mathcal U_\nu}(x)&=x^{q+1}-\nu,\\
 F_{K_e}(x)&=x^q-x-e/\omega,&
 F_{M_\delta}(x)&=x^{q-1}-\delta.
\end{aligned}
\end{equation}
In particular, all four polynomials have only simple roots.
\end{lemma}

\begin{theorem}\label{thm:fourth-grs}
For every integer $1\le k\le q-u-v-b-1$, there exists $\boldsymbol v=(v_\alpha)_{\alpha\in A}\in (\mathbb F_{q^2}^{*})^{N_4}$ such that $\GRS_k(A,\boldsymbol v)$ is a Hermitian self-orthogonal $[N_4,k,N_4-k+1]_{q^2}$ MDS code.
\end{theorem}

\begin{proof}
By Proposition~\ref{prop:intersection-formulas},
\[
 |L_C\cap K_E|=(u+1)v,
 \qquad
 |L_C\cap M_D|=ub,
 \qquad
 |K_E\cap M_D|=vb.
\]
The triple intersection has size $\tau_D$. Hence the inclusion--exclusion principle yields $|A|=N_4$. Since $u+v+b\le q-2$, the sum of the cardinalities of the three deleted sets is at most $q^2-q-b$. Thus $N_4\ge q+b>0$.

Define
\[
 P_C(x)=\prod_{c\in C}(x^q+x-c),
 \qquad
 P_E(x)=\prod_{e\in E}(\omega(x^q-x)-e),
\]
\[
 P_M(x)=\prod_{\delta\in D}(x^{q-1}-\delta).
\]
By Lemma~\ref{lem:basic-vanishing-polynomials}, we have $P_C=F_{L_C}$, $P_E=\omega^vF_{K_E}$, and $P_M=F_{M_D}$.

Since $0\in L_0$, we have $0\notin A$. In Lemma~\ref{lem:q-power-calculation}, take $X=\mathbb F_{q^2}$, $H=P_CP_E$, and $Q=P_M$. Then $P=P_CP_EP_M=HQ$, $c_0=\omega^v$, $b\ge1$, and $D\subseteq U$. For every $\alpha\in A$, the factors of the first two polynomials have the forms $\Tr(\alpha)-c$ and $\Tr(-\omega\alpha)-e$, respectively. Since $\alpha\notin L_C\cup K_E$, all these factors belong to $\mathbb F_q^{*}$, and hence $H(\alpha)\in\mathbb F_q^{*}$. Comparing degrees gives $\deg H=(u+v+1)q$.

For every positive integer $k$ with $k\le q-u-v-b-1$,
\[
\begin{aligned}
 &q^2+q-1-bq-\deg H-(q+1)k\\
 &\qquad=(q+1)(q-u-v-b-1-k)+(u+v+b)\ge0.
\end{aligned}
\]
The result therefore follows from Lemma~\ref{lem:q-power-calculation}.
\end{proof}

\end{document}
